\section*{अध्याय \ref{ch:g7:plate-tectonics} --- प्लेट विवर्तनिकी: चलती हुई पहेली}

\begin{solution}{exo:g7:plate-tectonics:1}
\omterm{def:g7:plate-tectonics:continental-drift}{महाद्वीपीय विस्थापन} यह विचार है कि \omterm{def:g3:continents-and-oceans:continent}{महाद्वीप} पृथ्वी की सतह पर धीरे-धीरे चलते
हैं, जिससे युगों के दौरान उनकी स्थितियाँ बदलती रही हैं। वेगनर के सुराग:
तटरेखाओं का मेल (दक्षिण अमेरिका और अफ़्रीका); आज \omterm{def:g2:water-on-earth:ocean}{महासागरों} से अलग \omterm{def:g3:continents-and-oceans:continent}{महाद्वीपों} पर
एक ही जीवाश्म; किसी \omterm{def:g2:water-on-earth:ocean}{महासागर} के दोनों ओर एक ही \omterm{def:g3:stones-pebbles-sand:rock}{चट्टानें} और पुरानी
\omterm{def:g4:world-physical-map:range}{पर्वत-श्रेणियाँ}; बेमेल जगहों पर प्राचीन \omterm{def:g5:weather-and-climate:climate}{जलवायु} के निशान (अंटार्कटिका में कोयला,
गरम, शुष्क अफ़्रीका में हिमनदों के निक्षेप)।
\end{solution}

\begin{solution}{exo:g7:plate-tectonics:2}
अपसारी: \omterm{def:g7:plate-tectonics:plate}{प्लेटें} एक-दूसरे से दूर जाती हैं। अभिसारी: वे एक-दूसरे की ओर आती
हैं। रूपांतर: वे अगल-बगल, एक-दूसरे के पास से सरकती हैं।
\end{solution}

\begin{solution}{exo:g7:plate-tectonics:3}
उदाहरण के लिए, प्रशांत \omterm{def:g2:water-on-earth:ocean}{महासागर} के चारों ओर (जापान, इंडोनेशिया, एंडीज़), और
भूमध्य सागर से हिमालय तक जाने वाली पट्टी। अफ़्रीका या ऑस्ट्रेलिया के बीच में
लगभग कोई बड़ा \omterm{def:g4:volcanoes-earthquakes-first-look:earthquake}{भूकंप} नहीं (रॉकी \omterm{def:g1:landscapes-around-us:landform}{पर्वतों} के पूर्व में उत्तर अमेरिका के, या
ब्राज़ील के बीच में भी नहीं)।
\end{solution}

\begin{solution}{exo:g7:plate-tectonics:4}
\omterm{def:g2:water-on-earth:salt-fresh}{मीठे पानी} का कोई जानवर खारे \omterm{def:g2:water-on-earth:ocean}{महासागर} के हज़ारों किलोमीटर पार नहीं कर सकता था।
उसका दोनों ओर मिलना यह बताता है कि जब \emph{मीसोसॉरस} रहता था, तब ये दोनों
जगहें ज़मीन से, एक ही \omterm{def:g3:continents-and-oceans:continent}{महाद्वीप} में, जुड़ी थीं।
\end{solution}

\begin{solution}{exo:g7:plate-tectonics:5}
कोयला गरम, नम, दलदली जंगलों के पौधों से बनता है: \omterm{def:g5:weather-and-climate:climate}{जलवायु} गरम और आर्द्र थी।
ऐसे जंगल आज \omterm{def:g3:continents-and-oceans:poles}{दक्षिणी ध्रुव} के पास नहीं उग सकते, इसलिए अंटार्कटिका कभी \omterm{def:g3:continents-and-oceans:equator}{भूमध्य
रेखा} के कहीं ज़्यादा पास रहा होगा: वह खिसका है।
\end{solution}

\begin{solution}{exo:g7:plate-tectonics:6}
मध्य-अटलांटिक कटक: अपसारी। एंडीज़: अभिसारी (एक महासागरीय \omterm{def:g7:plate-tectonics:plate}{प्लेट} एक \omterm{def:g3:continents-and-oceans:continent}{महाद्वीप} के
नीचे धँसती है)। सैन एंड्रियास भ्रंश: रूपांतर। हिमालय: अभिसारी (दो \omterm{def:g3:continents-and-oceans:continent}{महाद्वीप}
टकराते हैं)। पूर्वी अफ़्रीकी दरार \omterm{def:g1:landscapes-around-us:landform}{घाटी}: अपसारी (एक \omterm{def:g3:continents-and-oceans:continent}{महाद्वीप} फट रहा है)।
\end{solution}

\begin{solution}{exo:g7:plate-tectonics:7}
$2.5 \times 10 = 25$ किलोमीटर प्रति मिलियन वर्ष; 60 मिलियन वर्षों में:
$25 \times 60 = \qty{1500}{km}$।
\end{solution}

\begin{solution}{exo:g7:plate-tectonics:8}
$\qty{7}{m} = \qty{700}{cm}$, और $700 \div 9 \approx 78$: लगभग
\qty{78}{cm} प्रति वर्ष, जो \qty{2.5}{cm} प्रति वर्ष के औसत का लगभग 30 गुना है।
सतह पर \omterm{def:g7:plate-tectonics:plate}{प्लेटें} एकसार ढंग से अलग नहीं सरकतीं: खिंचाव वर्षों तक जमा होता रहता है
और छोटे-छोटे झटकों में निकलता है, जब दरारें एक ही बार में कई मीटर खुल जाती
हैं। लंबे समय में शांत वर्षों और झटकों का औसत कुछ सेंटीमीटर प्रति वर्ष
निकलता है।
\end{solution}

\begin{solution}{exo:g7:plate-tectonics:9}
वे यह नहीं बता सके कि कौन-सा बल \omterm{def:g3:continents-and-oceans:continent}{महाद्वीपों} को चलाता है, और उन्होंने जिस गति की
कल्पना की थी (\omterm{def:g3:continents-and-oceans:continent}{महाद्वीप} \omterm{def:g2:water-on-earth:ocean}{समुद्र} की तली को चीरते हुए आगे बढ़ते हैं) वह असंभव थी।
\omterm{def:g7:plate-tectonics:plate-tectonics}{प्लेट विवर्तनिकी} में \omterm{def:g3:continents-and-oceans:continent}{महाद्वीप} किसी चीज़ को चीरकर आगे नहीं बढ़ते: उन्हें, उनके
चारों ओर की समुद्री तली के साथ, \omterm{def:g7:plate-tectonics:plate}{प्लेटें} ढोती हैं, जो पृथ्वी के बाहरी कवच के
पूरे टुकड़ों की तरह चलती हैं।
\end{solution}

\begin{solution}{exo:g7:plate-tectonics:10}
सबसे ऊँची कटक पर, सबसे गहरी खाई पर। नई \omterm{def:g7:plate-tectonics:plate}{प्लेट} कटक पर बनती है, जहाँ \omterm{def:g6:inside-the-earth:layers}{प्रावार} से
गरम \omterm{def:g3:stones-pebbles-sand:rock}{चट्टान} ऊपर उठती है; महासागरीय \omterm{def:g7:plate-tectonics:plate}{प्लेट} खाई पर वापस \omterm{def:g6:inside-the-earth:layers}{प्रावार} में जाती है, जहाँ
वह \omterm{def:g3:continents-and-oceans:continent}{महाद्वीप} के नीचे धँसती है।
\end{solution}

\begin{solution}{exo:g7:plate-tectonics:11}
$14.2 \div 2.2 \approx 6.45$: लगभग 6 गुना तेज़।
किलोमीटर प्रति मिलियन वर्ष में चालें 142 और 22 हैं।
\qty{1000}{km} के लिए समय: पूर्वी प्रशांत उभार के लिए $1000 \div 142 \approx 7$
मिलियन वर्ष; उत्तर अटलांटिक कटक के लिए $1000 \div 22 \approx 45$ मिलियन
वर्ष।
\end{solution}

\begin{solution}{exo:g7:plate-tectonics:12}
(क) $6472 - 4578 = 1894$ \omterm{def:g4:volcanoes-earthquakes-first-look:earthquake}{भूकंप}। (ख) $4578 \div 7031 \approx 0.651$,
यानी 65\,\%। (ग) \omterm{def:g2:drawing-to-map:map}{मानचित्र} पर बनी रेखा खाई को दिखाती है, जहाँ महासागरीय \omterm{def:g7:plate-tectonics:plate}{प्लेट}
धँसना शुरू करती है। \omterm{def:g7:plate-tectonics:plate}{प्लेट} सीधी नीचे नहीं जाती: वह दूसरी \omterm{def:g7:plate-tectonics:plate}{प्लेट} के नीचे तिरछी
धँसती है, और उसके नीचे जितनी आगे जाती है, उतनी गहरी होती जाती है। \omterm{def:g4:volcanoes-earthquakes-first-look:earthquake}{भूकंप} पूरी
धँसती \omterm{def:g7:plate-tectonics:plate}{प्लेट} के साथ-साथ आते हैं, इसलिए उनमें से बहुत-से ऊपर वाली \omterm{def:g7:plate-tectonics:plate}{प्लेट} के नीचे,
खाई के एक ही ओर, रेखा से दसियों या सैकड़ों किलोमीटर दूर होते हैं। \omterm{def:g7:plate-tectonics:plate-boundary}{प्लेट सीमा}
एक क्षेत्र है, कोई पतली रेखा नहीं।
\end{solution}

\begin{solution}{pb:g7:plate-tectonics:1}
\textbf{1.} \emph{मीसोसॉरस}, \emph{साइनोग्नैथस} और \emph{ग्लोसॉप्टेरिस}
(\emph{लिस्ट्रोसॉरस} की पट्टी अफ़्रीका, भारत और अंटार्कटिका को जोड़ती है)।

\textbf{2.} वह \omterm{def:g2:water-on-earth:salt-fresh}{मीठे पानी} में रहता था: वह हज़ारों किलोमीटर चौड़ा खारा \omterm{def:g2:water-on-earth:ocean}{महासागर}
पार नहीं कर सकता था।

\textbf{3.} उत्तर-पूर्वी ब्राज़ील का उभार अफ़्रीका के पश्चिमी \omterm{def:g1:landscapes-around-us:water}{तट} पर गिनी की
खाड़ी के खोखले हिस्से में बैठता है।

\textbf{4.} एक \omterm{def:g7:plate-tectonics:plate-boundary}{अपसारी सीमा}, मध्य-अटलांटिक कटक: दक्षिण अमेरिकी \omterm{def:g7:plate-tectonics:plate}{प्लेट} पश्चिम की
ओर और अफ़्रीकी \omterm{def:g7:plate-tectonics:plate}{प्लेट} पूर्व की ओर, एक-दूसरे से दूर, चलती हैं।

\textbf{5.} दोनों \omterm{def:g1:landscapes-around-us:water}{तट} \omterm{def:g6:latitude-longitude:longitude}{प्रमुख याम्योत्तर} के विपरीत ओर हैं:
$43.2 + 14.4 = \ang{57.6}$।

\textbf{6.} $57.6 \times 102 = 5875.2$: लगभग \qty{5880}{km}।

\textbf{7.} \omterm{def:g3:continents-and-oceans:equator}{भूमध्य रेखा} से किसी \omterm{def:g3:continents-and-oceans:poles}{ध्रुव} की ओर जाने पर \omterm{def:g6:latitude-longitude:longitude}{याम्योत्तर} एक-दूसरे के पास
आते जाते हैं। \ang{23} दक्षिण वाली \omterm{def:g6:latitude-longitude:latitude}{अक्षांश रेखा} \omterm{def:g3:continents-and-oceans:equator}{भूमध्य रेखा} से छोटा वृत्त है,
इसलिए उसके 360 अंशों में से हर एक छोटा होता है।

\textbf{8.} लगभग आधी चौड़ाई: $5880 \div 2 = \qty{2940}{km}$।

\textbf{9.} $3.4 \times 10 = 34$ किलोमीटर प्रति मिलियन वर्ष।

\textbf{10.} $5880 \div 34 \approx 173$: लगभग 170 मिलियन वर्ष।

\textbf{11.} हाँ: 170 मिलियन वर्ष 100 और 200 मिलियन वर्षों के बीच है, और
\omterm{def:g7:plate-tectonics:supercontinent}{पैंजिया} के टूटने की शुरुआत, 225 से 200 मिलियन वर्ष पहले, के बाद आता है। \omterm{def:g2:drawing-to-map:map}{मानचित्र}
पर एक सीधा-सा हिसाब सही कोटि का
जवाब देता है।

\textbf{12.} उदाहरण के लिए: चाल हर समय एक जैसी नहीं रही; टूटने से पहले \omterm{def:g3:continents-and-oceans:continent}{महाद्वीप}
खिंचे और पतले हुए, इसलिए आज के \omterm{def:g1:landscapes-around-us:water}{तट} ठीक वह रेखा नहीं हैं जिसके साथ वे टूटकर अलग
हुए।
\end{solution}
